% !TeX root = ../tcc.tex
\chapter{Conclusões e limitações}
\label{cap:conclusoes}

Este trabalho reuniu as evidências sobre a avaliação de modelos de visão
computacional aplicados à identificação de manifestações cutâneas
associadas à hanseníase e, com base nas lacunas encontradas, planejou um
estudo de ciência de dados para avaliar dois classificadores. Este
capítulo apresenta as conclusões de cada etapa, as limitações do trabalho
e a continuidade prevista.

\section{Conclusões}
\label{sec:conclusoes}

A revisão estruturada incluiu 26 estudos, dos quais 17 são estudos
primários com resultados (Capítulo~\ref{cap:revisao-literatura}). As
evidências reunidas permitem as seguintes conclusões.

\begin{enumerate}
	\item \textbf{Técnicas.} Predominam as redes convolucionais, em geral
	      treinadas por transferência a partir de pesos pré-treinados, ao
	      lado de classificadores tradicionais aplicados a descritores de
	      textura ou a atributos extraídos por redes profundas
	      (Subseções~\ref{subsec:rq1-tecnicas} e~\ref{subsec:rq6-modelos}).
	\item \textbf{Avaliação.} A maior parte dos estudos avalia o desempenho
	      em uma única divisão da própria base: 14 dos 16 estudos que
	      desenvolvem modelos usam essa divisão, apenas 2 declaram separar os
	      dados por paciente, e só um testa o modelo em outra base
	      (Subseções~\ref{subsec:rq3-metodos} e~\ref{subsec:rq7-generalizacao}).
	      Sem separação por paciente, os indicadores podem estar inflados
	      por vazamento \cite{kapoor2023}.
	\item \textbf{Métricas.} A acurácia e a sensibilidade são as métricas
	      mais relatadas, enquanto a acurácia balanceada e a calibração dos
	      escores não aparecem nos estudos com dados
	      (Subseção~\ref{subsec:rq4-metricas}). Os estudos que isolam a
	      hanseníase relatam sensibilidade bem inferior aos indicadores
	      globais.
	\item \textbf{Equidade.} Nenhum dos 17 estudos com dados reporta
	      desempenho por tom de pele (Subseção~\ref{subsec:rq5-rq6-tom-pele}).
	\item \textbf{Comparabilidade.} Bases, classes e procedimentos diferem
	      entre os estudos, e há sobreposição de origem entre conjuntos
	      públicos, o que impede comparar numericamente os resultados
	      relatados.
\end{enumerate}

O planejamento do Capítulo~\ref{cap:metodologia-experimental} responde a
essas lacunas. Ele define a fotografia como unidade de avaliação, mantém
fotografias relacionadas no mesmo grupo em todas as partições, seleciona
os hiperparâmetros por validação cruzada estratificada por grupos, adota a
acurácia balanceada como medida principal, relata intervalos por
\textit{bootstrap} de grupos e a diferença pareada entre as duas redes, e
inclui linhas de base que verificam se atributos da fonte bastam para
prever a classe. Os atributos essenciais do padrão \textit{Data Science}
da ACM SIGSOFT estão relacionados às partes do trabalho na
Tabela~\ref{tab:atributos-data-science}.

\section{Limitações}
\label{sec:limitacoes}

\subsection{Limitações da revisão}

\begin{itemize}
	\item \textbf{Estratégia de busca.} A \textit{string} booleana retornou
	      10.818 artigos na Scopus mesmo após diversos ajustes, e a busca
	      principal foi feita com o Consensus e o Claude
	      (Seção~\ref{subsec:picoc-string-busca}). Essas ferramentas
	      recuperam publicações por similaridade e não expõem integralmente
	      seus critérios internos, o que limita a reprodutibilidade da busca
	      e pode ter deixado estudos relevantes de fora
	      (Apêndice~\ref{apend:registro-busca}).
	\item \textbf{Triagem.} As sugestões do Jev Reviewer serviram de apoio,
	      mas divergiram das decisões em boa parte dos incluídos, e a
	      documentação não permite afirmar que os dois pesquisadores
	      avaliaram todos os registros de forma independente
	      (Seção~\ref{sec:execucao-busca-selecao}).
	\item \textbf{Critérios de seleção.} A formulação preservada do critério
	      CI03 está truncada, e o texto integral ainda deve ser confirmado
	      (Apêndice~\ref{apend:registro-busca}).
	\item \textbf{Extração.} A extração depende do que os artigos relatam.
	      Dados não localizados, como o número de pacientes ou a unidade da
	      divisão, foram registrados como tal, sem que se possa afirmar que
	      inexistam.
	\item \textbf{Síntese.} A heterogeneidade de bases, tarefas, classes e
	      procedimentos permitiu apenas uma síntese descritiva, sem
	      metanálise.
\end{itemize}

\subsection{Limitações do planejamento}

\begin{itemize}
	\item \textbf{Ausência de resultados.} Este trabalho apresenta o plano do
	      estudo; os resultados dependem da execução prevista para a
	      Monografia~2.
	\item \textbf{Dados.} O estudo usará uma única coleção pública, com
	      rótulos sem confirmação clínica, sem identificação de pacientes e
	      com fontes de origem comum a conjuntos usados na literatura. O
	      teste será interno a essa coleção.
	\item \textbf{Escopo da tarefa.} A tarefa é binária e não distingue as
	      formas clínicas da hanseníase nem as condições da classe negativa.
	\item \textbf{Análises exploratórias.} As análises por tom de pele, de
	      calibração e de predição seletiva dependem do tamanho da amostra
	      curada e não sustentam conclusões confirmatórias.
\end{itemize}

\section{Continuidade}
\label{sec:continuidade}

Na Monografia~2, estão previstas a curadoria do 32 CSD conforme o
protocolo, a geração e o registro das partições, o treinamento e a seleção
das duas redes, a definição dos limiares, a abertura do teste reservado e
a análise dos resultados frente às questões do GQM. Em seguida, os
resultados serão discutidos em relação às evidências da revisão, com
atenção às ameaças à validade descritas na
Seção~\ref{sec:validade-avaliacao-empirica}. Se o acesso ao AI4Leprosy for
concedido, a mesma avaliação será repetida nessa coleção, sem alterar o
protocolo, o que permitirá testar os classificadores em fotografias de
outra origem.
